\section{Descarte’s Nightmare}

\begin{quotex}
The battle differentiates, selects, creates hierarchy; especially when -- to use the traditional terms -- it is not the lesser battle, but the greater battle; not the battle of man against man, or against the world, but the battle of the supernatural element of man against everything in him that is nature, sensation, materiality, ferment, delusion of grandeur; against the chaos that is within him, before it is outside him.
\flright{\textsc{Julius Evola}}
\end{quotex}


\begin{quotex}
The rule and criterion of truth is to have made it. Hence the clear and distinct idea of the mind not only cannot be the criterion of other truths, but it cannot be the criterion of that of the mind itself; for while the mind apprehends itself, it does not make itself, and because it does not make itself it is ignorant of the form or mode by which it apprehends itself.
\flright{\textsc{Giambattista Vico} (commenting on Descartes)}
\end{quotex}


The fundamental difference between the traditional and modern world views is the knowledge of a spiritual level of reality that transcends our ordinary experience. This spiritual level is the source of all manifestation and is known, not through empirical evidence or a process of reasoning, but rather through a direct ``knowing'', called Intuition or gnosis. This gnosis is not irrational, but is supra-rational, that is, above reason, though not contrary to reason.

The traditional teaching is that man is a tripartite being, that is, a Unity of Spirit, soul, and body. These states of being are hierarchically arranged, with the lower state subordinated to the higher states.

The modernist denies the existence of a spiritual state of being, and thus, of any knowledge that transcends the senses or the reasoning faculty. Science ignores it completely by its very methodology, and regards any claim to the contrary as delusional. The major exoteric religions accept the existence of the spiritual realm, but insist it can be known only by faith, with a caste of priests or ministers keeping guard over it.

Now, the difference between the Traditional and modern world view must not be understood in a temporal sense, but rather ontologically. ``Tradition'' indicates a transcendental source that impinges on the empirical at every time and place. ``Modern'' indicates the denial of such transcendence, viewing the empirical as self-contained. Thus, the Traditional and modern world views are independent of time and place, and have always been with us, with one predominating over the other. Other times, these two spirits are found battling each other for influence.\footnote{For the remainder of the essay, please go to Chapter \ref{sec:DescartesNightmare} in this book.}


\flrightit{Posted on 2010-03-14 by Cologero}

